% TOP_PROBLEM: {"id": 3046, "title": "Beneš Conjecture", "category_id": 2, "category": {"id": 2, "name": "combinatorics", "display_name": "Combinatorics", "description": "Counting problems, graph theory, discrete structures.", "slug": "combinatorics", "order_index": 2, "created_at": "2026-07-31T15:26:25.671Z"}, "status": "open"}
% FIRST_POSTED: null
% ENTRY_KIND: statement_only
% Imported from: batch14.tex; UnsolvedMath ID 3046
\documentclass[11pt]{article}
\usepackage[margin=25mm]{geometry}
\usepackage{fontspec}
\setmainfont{Latin Modern Roman}
\usepackage{amsmath,amssymb,mathtools}
\usepackage{unicode-math}
\setmathfont{Latin Modern Math}
\usepackage{xurl,hyperref}
\hypersetup{colorlinks=true,urlcolor=blue}
\setcounter{secnumdepth}{0}
\setlength{\parindent}{0pt}
\setlength{\parskip}{5pt}
\setlength{\emergencystretch}{3em}
\newcommand{\sref}[2]{\hyperlink{source:#1:#2}{[S#2]}}
\newcommand{\cataloguescope}[1]{\paragraph{Recorded target.}#1}
\begin{document}
% UnsolvedMath ID 3046; source catalogue code OPG-37181
\setcounter{section}{3045}
\section{Beneš Conjecture}\label{3046}
{\small\sffamily UnsolvedMath ID 3046 · Combinatorics}
\subsection{Definitions and mathematical statement}

\paragraph{Shared notation.}$\mathbb N=\{1,2,\ldots\}$, and $\mathbb Z,\mathbb Q,\mathbb R,\mathbb C$ denote the integers, rationals, reals and complex numbers. Any other conventions are specified locally.


Let permutations of a nonempty finite set $E$ act on the right, with $x(\alpha\beta)=(x\alpha)\beta$. A partition $h$ of $E$ is a family of nonempty pairwise disjoint blocks whose union is $E$. Its blockwise stabilizer is
$S(h)=\{\pi\in\mathfrak S(E):B\pi=B\text{ for every }B\in h\}$, where $B\pi=\{x\pi:x\in B\}$.
For a permutation $\varphi$, write $\varphi(h)=\{B\varphi:B\in h\}$ for the image partition. The meet $h\wedge h'$ consists of all nonempty intersections of their blocks, and $0_E=\{\{x\}:x\in E\}$ is the singleton partition. A partition is uniform if all blocks have the same cardinality. A subset $P\subset\mathfrak S(E)$ is transitive if for every $x,y\in E$ some $\pi\in P$ sends $x$ to $y$. Products and powers are setwise products of exactly the indicated number of factors.
\paragraph{Statement recorded in the supplied source.}
Find a sufficient condition for a sequence of partitions $h_1,\ldots,h_\ell$ to be complete, meaning $S(h_1)\cdots S(h_\ell)=\mathfrak S(E)$. In particular, characterize useful sufficient conditions for the orbit sequence $h,\delta(h),\ldots,\delta^{\ell-1}(h)$ of a partition under a permutation $\delta$.

Resolve the recorded Beneš implication: for every uniform partition $u$, every $\varphi\in\mathfrak S(E)$ with $u\wedge\varphi(u)=0_E$, and every integer $n\geq2$, does transitivity of $(\varphi S(u))^n$ imply
\[
 (\varphi S(u))^{2n-1}=\mathfrak S(E)?
\] \sref{3046}{2}
\subsection{Short English statement}
Find conditions making a product of partition stabilizers equal the whole symmetric group, and test the Beneš transitivity criterion for an iterated partition.
\subsection{Sources}
\begin{description}
\item[\hypertarget{source:3046:1}{S1}] ULAM.AI and the UnsolvedMath Contributors, \emph{UnsolvedMath: A Curated Collection of Open Mathematics Problems}, record 3046 (OPG-37181). CC BY 4.0. \url{https://huggingface.co/datasets/ulamai/UnsolvedMath}.
\item[\hypertarget{source:3046:2}{S2}] Przemek Chojecki, \emph{Beneš and Shuffle-Exchange Counterexamples}, arXiv:2607.15296 (2026), §2, Corollary 3, and §4, Theorem 19. Corollary 3 refutes the Beneš partition-stabilizer implication; Theorem 19 gives $d(k,3)=6$ for $k\geq3$ and refutes the universal shuffle-exchange formula. \url{https://arxiv.org/abs/2607.15296}.
\end{description}
\label{end:3046}
\end{document}
